
\chapter{Metodologia proposta}

A pesquisa terá caráter experimental e quantitativo, com execução de experimentos computacionais controlados para comparar estratégias de aprendizado contrastivo na predição de propriedades de materiais bidimensionais. O trabalho será desenvolvido em etapas: preparação dos dados da base C2DB, construção das representações em grafos, treinamento dos modelos e avaliação em tarefas de regressão. Os métodos CrystalCLR, de Koker et al.~\cite{koker2022graph}, e SMICLR, de Pinheiro et al.~\cite{Pinheiro2022}, serão utilizados como referência para a adaptação das técnicas de aprendizado contrastivo ao domínio químico.
Na implementação, a CGCNN será adotada como arquitetura principal do codificador, por ser compatível com grafos cristalinos e permitir a comparação direta com o CrystalCLR. A mesma arquitetura e as mesmas partições dos dados serão usadas nas configurações contrastivas e no modelo de referência supervisionado, de modo que a comparação permita avaliar a contribuição do pré-treinamento contrastivo. A Figura \ref{fig:arquitetura_proposta} apresenta o fluxo geral do trabalho.


\begin{figure}[htbp]
    \centering
    \caption{Fluxo conceitual da arquitetura proposta para predição de propriedades de materiais bidimensionais.}
    \label{fig:arquitetura_proposta}
    \resizebox{\textwidth}{!}  {%
    \begin{tikzpicture}[
        font=\small,
        node distance=1.2cm and 1.3cm,
        every text node part/.style={align=center},
        % Estilos de Blocos
        boxInput/.style={
            rectangle, draw=mainblue, fill=lightblue, thick, rounded corners=5pt,
            minimum width=3.2cm, minimum height=1.1cm, font=\sffamily
        },
        boxProc/.style={
            rectangle, draw=accentcyan, fill=lightcyan, thick, rounded corners=5pt,
            minimum width=3.0cm, minimum height=1.1cm, font=\sffamily
        },
        boxModel/.style={
            rectangle, draw=accentpurple, fill=lightpurple, thick, rounded corners=5pt,
            minimum width=3.8cm, minimum height=1.2cm, font=\bfseries\sffamily
        },
        boxLoss/.style={
            rectangle, draw=accentgreen, fill=lightgreen, thick, rounded corners=5pt,
            minimum width=3.6cm, minimum height=1.1cm, font=\sffamily
        },
        boxOut/.style={
            rectangle, draw=mainblue, fill=lightblue, thick, rounded corners=5pt,
            minimum width=2.8cm, minimum height=1.1cm, font=\sffamily
        },
        arrow/.style={
            -{Stealth[scale=1.0]}, thick, color=darkgray
        },
        dashedArrow/.style={
            -{Stealth[scale=1.0]}, thick, color=accentpurple, dashed
        },
        groupTitle/.style={
            font=\bfseries\sffamily, color=darkgray
        }
    ]

    % --- FASE 1: Pré-Treinamento Contrastivo ---
    \node (c2db) [boxInput] {\textbf{Banco de Dados C2DB}\\Estruturas 2D};
    \node (graph) [boxProc, below=of c2db] {\textbf{Grafo Cristalino}\\$G = (V, E)$\\(Átomos e Arestas)};

    % Estratégias (Auto-supervisionada e Supervisionada)
    \node (aug) [boxProc, above right=0.4cm and 1.0cm of graph] {\textbf{Ramo Auto-Supervisionado}\\Aumentações Estruturais ($\tilde{G}_i, \tilde{G}_j$)};
    \node (sup) [boxProc, below right=0.4cm and 1.4cm of graph] {\textbf{Ramo Supervisionado}\\Pares por Rótulo ($E_g, \Delta H$)};

    % Codificadores GNN
    \node (encoder1) [boxModel, right=1.0cm of aug] {\textbf{Codificador GNN } $f_\theta$};
    \node (encoder2) [boxModel, right=1.85cm of sup] {\textbf{Codificador GNN } $f_\theta$};

    % Cabeças de Projeção
    \node (proj1) [boxProc, right=0.9cm of encoder1] {\textbf{Projeção } $g_\phi$\\Vetor $\mathbf{z}_i$};
    \node (proj2) [boxProc, right=0.9cm of encoder2] {\textbf{Projeção } $g_\phi$\\Vetor $\mathbf{z}_j$};

    % Função de Perda Contrastiva
    \node (loss) [boxLoss, right=1.0cm of proj1, yshift=-1.5cm] {\textbf{Perda Contrastiva} $\mathcal{L}_{\text{cont}}$\\[2pt]\scriptsize $\min d_{\cos}(\mathbf{z}_i, \mathbf{z}_j)$ (Positivos)\\\scriptsize $\max d_{\cos}(\mathbf{z}_i, \mathbf{z}_k)$ (Negativos)};

    % Conexões Fase 1
    \draw [arrow] (c2db) -- (graph);
    \draw [arrow] (graph.east) -- ++(0.4,0) |- (aug.west);
    \draw [arrow] (graph.east) -- ++(0.4,0) |- (sup.west);
    \draw [arrow] (aug) -- (encoder1);
    \draw [arrow] (sup) -- (encoder2);
    \draw [arrow] (encoder1) -- (proj1);
    \draw [arrow] (encoder2) -- (proj2);
    \draw [arrow] (proj1.east) -- ++(0.3,0) |- (loss.west);
    \draw [arrow] (proj2.east) -- ++(0.3,0) |- (loss.west);

    % Indicação de Pesos Compartilhados
    \draw [<->, dashed, color=accentpurple, thick] (encoder1) -- node[left, font=\scriptsize\sffamily, align=right] {\textbf{Pesos}\\\textbf{Compartilhados}} (encoder2);

    % Container Fase 1
    \begin{scope}[on background layer]
        \node (bg1) [
            draw=bordergray, fill=black!2, thick, rounded corners=8pt,
            fit=(c2db) (graph) (aug) (sup) (encoder1) (encoder2) (proj1) (proj2) (loss),
            inner sep=20pt
        ] {};
        \node [anchor=north west, groupTitle] at (bg1.north west) {FASE 1: Pré-Treinamento Contrastivo (Aprendizado de Representação)};
    \end{scope}

    % --- FASE 2: Ajuste Fino e Regressão Downstream ---
    \node (testData) [boxInput, below=4.0cm of graph] {\textbf{Estrutura 2D Alvo}\\$G \in \text{C2DB}$};
    \node (ftEncoder) [boxModel, right=2.0cm of testData] {\textbf{Codificador GNN } $f_\theta$\\(Pesos Pré-treinados)};
    \node (emb) [boxProc, right=1.5cm of ftEncoder] {\textbf{Vetor Latente}\\$\mathbf{h}\in\mathbb{R}^d$};
    \node (mlp) [boxProc, right=1.5cm of emb] {\textbf{Regressor MLP}\\Camadas Densas};
    \node (preds) [boxOut, right=1.5cm of mlp] {\textbf{Propriedades Alvo}\\[3pt] \small Band Gap ($E_g$)\\[2pt] \small Energia Formação ($\Delta H$)};

    % Conexões Fase 2
    \draw [arrow] (testData) -- (ftEncoder);
    \draw [arrow] (ftEncoder) -- (emb);
    \draw [arrow] (emb) -- (mlp);
    \draw [arrow] (mlp) -- (preds);

    % Transferência de Conhecimento (Fase 1 para Fase 2)
    \draw [dashedArrow] (encoder2.south) 
        -- ([yshift=0.45cm] encoder2.south |- ftEncoder.north) 
        -| node[pos=0.25, above, font=\scriptsize\sffamily, align=center] {\textbf{Transferência de Pesos $f_\theta$}} (ftEncoder.north);

    % Container Fase 2
    \begin{scope}[on background layer]
        \node (bg2) [
            draw=bordergray, fill=black!2, thick, rounded corners=8pt,
            fit=(testData) (ftEncoder) (emb) (mlp) (preds),
            inner sep=20pt
        ] {};
        \node [anchor=north west, groupTitle] at (bg2.north west) {FASE 2: Ajuste Fino (\textit{Downstream Task})};
    \end{scope}

    \end{tikzpicture}%
    }
    % \fonte{Elaborado pelo autor (2026).}
\end{figure}


\section{Procedimentos metodológicos}

\begin{enumerate}
    \item \textbf{Seleção e preparação dos dados:} serão obtidas da C2DB~\cite{haastrup2018computational,gantenbein2021c2db} as estruturas cristalinas e os valores disponíveis para o \textit{band gap} e a energia de formação. Serão verificadas a presença de valores ausentes, unidades, convenções de cálculo e estruturas duplicadas. Para cada propriedade, serão mantidas apenas as amostras com rótulos válidos e compatíveis; se houver mais de um nível de cálculo para uma propriedade, será selecionada uma única convenção para os experimentos comparáveis.

    \item \textbf{Construção dos grafos e partições:} cada estrutura será representada por um grafo periódico, com átomos como nós e vizinhanças atômicas como arestas, codificando-se as espécies químicas e as distâncias interatômicas. Após a remoção de duplicatas, as estruturas serão separadas em conjuntos de treinamento, validação e teste na proporção inicial de 70\%, 15\% e 15\%, respectivamente. Estruturas duplicadas ou equivalentes não serão distribuídas entre partições diferentes, para reduzir o risco de vazamento de dados. Caso o número de amostras válidas de uma propriedade não permita essa divisão, as proporções serão ajustadas e registradas antes do treinamento. A mesma divisão será utilizada em todas as comparações.

    \item \textbf{Definição e validação das aumentações:} serão implementadas duas visões de cada estrutura de treinamento, necessárias ao aprendizado contrastivo. Serão avaliadas transformações geométricas fisicamente plausíveis, como pequenas perturbações das posições atômicas e deformações no plano da monocamada. Os parâmetros das transformações serão limitados e verificados para evitar estruturas inválidas ou alterações que mudem a identidade química do material. As aumentações serão aplicadas somente aos dados de treinamento; validação e teste permanecerão sem alterações.

    \item \textbf{Treinamento dos modelos:} serão comparadas três configurações: (a) CGCNN treinada diretamente para regressão, como referência supervisionada; (b) pré-treinamento auto-supervisionado com a perda NT-Xent, usando duas visões aumentadas da mesma estrutura como par positivo e as demais amostras do lote como negativos; e (c) pré-treinamento supervisionado por contraste, usando a perda SupCon. Nesta última configuração, os valores contínuos de cada propriedade serão divididos em faixas definidas pelos quantis do conjunto de treinamento; materiais na mesma faixa formarão pares positivos e materiais em faixas distintas, pares negativos. Será treinado um modelo contrastivo supervisionado separado para cada propriedade. Após o pré-treinamento, o codificador será ajustado junto a uma camada de regressão usando os rótulos de treinamento. Os hiperparâmetros serão escolhidos com base no conjunto de validação, sem utilizar os resultados do teste.

    \item \textbf{Avaliação e análise:} os codificadores treinados serão avaliados em tarefas de regressão do \textit{band gap} e da energia de formação. O desempenho será medido por erro absoluto médio (MAE), raiz do erro quadrático médio (RMSE) e coeficiente de determinação ($R^2$), apresentados separadamente para cada alvo. Para investigar o efeito da escassez de rótulos, o ajuste fino será repetido utilizando 10\%, 25\%, 50\% e 100\% dos rótulos disponíveis no conjunto de treinamento; validação e teste permanecerão inalterados. Cada experimento será repetido com três sementes aleatórias, reportando-se a média e o desvio-padrão das métricas. A análise final discutirá o desempenho, a variabilidade e as limitações observadas, sem pressupor previamente que o método contrastivo será superior.

    \item \textbf{Registro e redação:} serão documentados as versões dos dados e das bibliotecas, as configurações dos experimentos, as sementes aleatórias e os resultados obtidos. Os resultados e as decisões metodológicas serão incorporados progressivamente ao texto, que será revisado ao final do período.
\end{enumerate}

\section{Atividades e Cronograma}

\subsection{Lista de atividades}

\begin{enumerate}
    \item Concluir o levantamento bibliográfico e fixar o protocolo experimental, incluindo arquitetura principal, propriedades-alvo, critérios de inclusão e métricas.
    \item Obter e caracterizar os dados da C2DB; tratar valores ausentes, unidades, duplicatas e a consistência dos rótulos.
    \item Implementar a conversão das estruturas em grafos periódicos, definir as partições de treinamento, validação e teste e verificar a ausência de sobreposição entre elas.
    \item Implementar e validar as aumentações estruturais e treinar o modelo de referência supervisionado.
    \item Implementar e treinar as configurações auto-supervisionada e supervisionada por contraste, ajustando hiperparâmetros com o conjunto de validação.
    \item Executar as tarefas de regressão e os experimentos com diferentes quantidades de rótulos; calcular e analisar as métricas.
    \item Consolidar os resultados, revisar a metodologia e redigir a versão final do trabalho.
\end{enumerate}

\subsection{Cronograma}

O cronograma proposto para o período de janeiro a junho de 2027 é apresentado na Tabela \ref{tab:cronograma}. As atividades de redação e registro dos experimentos ocorrerão ao longo de todo o período.

\begin{table}[htbp]
    \centering
    \caption{Cronograma proposto para o desenvolvimento do trabalho.}
    \label{tab:cronograma}
    \small
    \begin{tabular}{|p{7.1cm}|c|c|c|c|c|c|}
        \hline
        \textbf{Atividade} & \textbf{Jan.} & \textbf{Fev.} & \textbf{Mar.} & \textbf{Abr.} & \textbf{Mai.} & \textbf{Jun.} \\
        \hline
        Revisão bibliográfica e definição do protocolo experimental & X & X & & & & \\
        \hline
        Obtenção, análise e limpeza dos dados da C2DB & X & X & & & & \\
        \hline
        Conversão das estruturas em grafos e definição das partições & & X & X & & & \\
        \hline
        Implementação e validação das aumentações; modelo supervisionado de referência & & & X & X & & \\
        \hline
        Implementação e treinamento dos modelos contrastivos & & & & X & X & \\
        \hline
        Experimentos downstream, métricas e análise comparativa & & & & & X & X \\
        \hline
        Consolidação dos resultados, redação e revisão final & & & & X & X & X \\
        \hline
    \end{tabular}
\end{table}

